% !TeX root = ../main.tex
\chapter{Subsystems}
\label{ch:subsystems}

\section{Decomposition}

\umlfigure{subsystems}{Subsystems grouped by the project that contains them.
Dashed boundaries are projects; solid boxes are subsystems}{fig:subsystems}

Figure~\ref{fig:subsystems} shows six subsystems distributed across three
projects and the two external systems they talk to. The dashed boundaries are
the projects of the solution; the solid boxes inside them are the subsystems;
the grey boxes are systems we do not own.

\begin{longtable}{L{0.18\textwidth}L{0.21\textwidth}L{0.30\textwidth}L{0.17\textwidth}}
	\caption{Subsystems and their responsibilities}\label{tab:subsystems}\\
	\toprule
	\textbf{Subsystem} & \textbf{Project} & \textbf{Responsibility} & \textbf{Requirements} \\
	\midrule
	\endfirsthead
	\toprule
	\textbf{Subsystem} & \textbf{Project} & \textbf{Responsibility} & \textbf{Requirements} \\
	\midrule
	\endhead

	Repository Client & \code{Pentateuch.Genesys} & Authenticates against a
	GENESYS repository, opens a named project, and reads the schema that
	project actually defines. & \req{CM-1}, \req{CM-2} \\
	\addlinespace
	Model Mapper & \code{Pentateuch.Genesys} & Converts between GENESYS
	entities and the abstract syntax, in both directions. Creates entities and
	relationships only. & \req{SA-2}, \req{AS-2}, \req{CA-2}, \req{CM-5} \\
	\addlinespace
	Code Extractor & \code{Pentateuch.Translator} & Opens a solution, resolves
	symbols through the semantic model, builds a control flow graph per method
	reached. & \req{CA-1} \\
	\addlinespace
	Code Generator & \code{Pentateuch.Translator} & Emits C\# or Visual Basic
	implementing the control flow of an activity diagram, leaving method
	bodies as prototypes. & \req{AC-1} to \req{AC-4} \\
	\addlinespace
	Abstraction & \code{Pentateuch.Translator} & Discards constructs below the
	level a systems engineer reads. The one subsystem that exists for a
	non-functional requirement rather than a functional one. & \req{NFR-8} \\
	\addlinespace
	Abstract Syntax & \code{Pentateuch.Absyn} & The form both sides exchange.
	Belongs to neither the language nor the model. & The pivot of
	Section~\ref{sec:canonical} \\
	\bottomrule
\end{longtable}

\section{Why the Boundaries Fall Where They Do}

The project boundaries were not chosen for tidiness. Two constraints fixed them.

\textbf{C-1} is the first. The GENESYS Client API assemblies target .NET
Framework 4.8, and code calling them must target the same framework. The Roslyn
work wants a current runtime. No single project can do both, so the split across
runtimes is imposed rather than preferred, and the subsystems that touch each
vendor API are grouped by which runtime that API demands.

\textbf{C-3} is the second. GENESYS loads its own copy of Roslyn in order to
compile the scripts it supports, and two versions of the same assembly cannot
coexist in one process. The analysis therefore cannot run inside GENESYS even if
the framework problem were solved. Taken together with \req{NFR-2}, a failure in
the tool must not terminate GENESYS, the two processes are settled.

\code{Pentateuch.Absyn} targets .NET Standard 2.0 because that is the only
target both of the other runtimes can reference. Its existence is a consequence
of the same constraint that created the split.

The arrangement has a second benefit. Two of the three projects reference no
vendor assembly, so they build and run on a
machine with no GENESYS installation and no licence activation. That is
\req{NFR-6}, verified by test V-12, and it is what keeps the limited number of
licence activations from becoming a limit on how many people can work.

\section{The Canonical Form}
\label{sec:canonical}

\umlfigure{canonical-model}{The abstract syntax as the canonical form between
the two outer representations}{fig:canonical}

Figure~\ref{fig:canonical} shows why the middle of the pipeline is a
representation in its own right rather than an incidental data structure.

A translation is a chain of three representations, each of which knows how to
translate to the next and exposes a \code{translate()} operation. The source
representation is Roslyn's view of the program; the model representation is the
Client API's view of the repository; the abstract syntax sits between them and
belongs to neither.

This reduces how many translators are needed. If each format translated
directly to every other, supporting $N$ formats would need on the order of
$2N$ translators. Routing
everything through one canonical form makes it $N$. Concretely, Visual Basic
becomes a second translation into the same abstract syntax rather than a change
to anything downstream, which is the reason \req{AC-1} can name two target
languages without doubling the work. Excluding C from scope keeps both supported
languages on one compiler platform, so the two spokes share a front end as well.

\section{How the Subsystems Work Together}

A translation is one pass along the pipeline, with the subsystems entering in a
fixed order. The order is fixed by the design rather than by convention, which
is what makes \req{NFR-4}, identical output from identical input, achievable at
all.

\subsection{Code to Activity Diagram}

The Code Extractor opens the solution, resolves its project references and symbols through the compiler's semantic model (so each call is matched to the method it actually invokes, not just a name), and builds a control flow graph for every method it reaches. It works from a queue seeded with the engineer's chosen entry point and limited to their designated depth, which is
\req{CA-3}.

The Abstraction filter receives the completed graph and applies the significance
rules to it. It runs after the walk rather than during it, because filtering
during traversal would make the discarded material unavailable for the counts
reported at the end.

The Model Mapper receives abstract syntax and produces the entities and
relationships that correspond to what survived. It asks the Repository Client
for the schema first, so that it addresses classes by what the schema reports
rather than by assumed names, which is \req{CM-2}.

The Repository Client holds the proposed change set and writes nothing until the
user accepts it. On acceptance it applies the whole set; on cancellation it
applies none of it. That is \req{CM-3} and \req{NFR-3}.

\subsection{Activity Diagram to Code}

The same subsystems in the reverse order. The Repository Client reads, the Model
Mapper produces abstract syntax from entities, the Abstraction filter applies
the same significance rules, and the Code Generator emits source with method
bodies left as prototypes for an engineer to complete. The review step is the
same one, applied to the source repository instead of the model.

\subsection{What Crosses Each Boundary}

\begin{longtable}{L{0.38\textwidth}L{0.54\textwidth}}
	\caption{Interfaces between subsystems}\label{tab:interfaces}\\
	\toprule
	\textbf{Boundary} & \textbf{What crosses it} \\
	\midrule
	\endfirsthead
	\toprule
	\textbf{Boundary} & \textbf{What crosses it} \\
	\midrule
	\endhead

	Source Repository to Code Extractor & Source files on disk, opened as a
	solution so that references resolve. \\
	\addlinespace
	Code Extractor to Abstraction & Syntax trees and control flow graphs,
	still carrying every construct the source contained. \\
	\addlinespace
	Abstraction to Model Mapper & Abstract syntax, carrying only what survived
	the significance rules. This is the process boundary, and the open
	question of risk \textbf{R-2}. \\
	\addlinespace
	Model Mapper to Repository Client & Entities and relationships, named
	according to the schema read at run time. \\
	\addlinespace
	Repository Client to GENESYS & Client API calls, under a licence, against
	one open project. \\
	\bottomrule
\end{longtable}

